% Generated by article_paper_simulations.R; requires booktabs and graphicx.
\begin{table}[p]\centering
\caption{VAR lag selection by BIC.}
\label{tab:sim-var-lag-selection}
\small\setlength{\tabcolsep}{4pt}\renewcommand{\arraystretch}{1.15}
\resizebox{\linewidth}{!}{%
\begin{tabular}{llrrrrrr}\toprule
Design & $n$ & Returned & Max. lag & Mean lag & Median lag & At max. (\%) & Capped \\\midrule
1 & 100 & 1000 & 10 & 1.000 & 1.000 & 0.000 & 0 \\
1 & 200 & 1000 & 13 & 1.001 & 1.000 & 0.000 & 0 \\
1 & 400 & 1000 & 15 & 1.000 & 1.000 & 0.000 & 0 \\
2 & 100 & 1000 & 10 & 2.238 & 2.000 & 0.000 & 0 \\
2 & 200 & 1000 & 13 & 3.048 & 3.000 & 0.000 & 0 \\
2 & 400 & 1000 & 15 & 4.621 & 5.000 & 0.000 & 0 \\
3 & 100 & 1000 & 10 & 1.022 & 1.000 & 0.000 & 0 \\
3 & 200 & 1000 & 13 & 1.062 & 1.000 & 0.000 & 0 \\
3 & 400 & 1000 & 15 & 1.263 & 1.000 & 0.000 & 0 \\
4 & 100 & 1000 & 10 & 2.008 & 2.000 & 0.000 & 0 \\
4 & 200 & 1000 & 13 & 2.026 & 2.000 & 0.000 & 0 \\
4 & 400 & 1000 & 15 & 2.193 & 2.000 & 0.000 & 0 \\
\bottomrule\end{tabular}}
\par\medskip
\begin{minipage}{\linewidth}\footnotesize The mean, median and boundary frequency use returned fits. BIC compares candidate orders on a common sample; the selected order is refitted by OLS using all available observations. Capped counts bounds reduced by the residual degrees-of-freedom safeguard. The selected VAR lag is an approximation order and is not compared with the DGP's VARMA AR order. Full frequencies are in data/var\_lag\_frequencies.csv.\end{minipage}
\end{table}
